\chapter{Implementation Details}

This appendix records implementation-level details that are useful for replication but are
not central to the data-collection rationale discussed in Chapter 3.

\section{Frontend Implementation}

The frontend is a React 18.2 SPA built with Vite. It uses a role-based component structure
(participant vs. admin) and React Context for authentication and WebSocket session state.

\subsection{Project Structure}
\begin{verbatim}
frontend/
|-- src/
|   |-- components/
|   |   |-- PrivateRoute.jsx
|   |   |-- admin/
|   |   |   |-- AdminDashboard.jsx
|   |   |   |-- DebateManagement.jsx
|   |   |   |-- MatchmakingDashboard.jsx
|   |   |   +-- UserManagement.jsx
|   |   |-- auth/
|   |   |   |-- GuestResumeForm.jsx
|   |   |   +-- Login.jsx
|   |   +-- participant/
|   |       |-- DebateInterface.jsx
|   |       |-- DebateRoom.jsx
|   |       |-- JoinDebate.jsx
|   |       |-- ParticipantDashboard.jsx
|   |       +-- PostDebateSurveyModal.jsx
|   |-- context/
|   |   |-- AuthContext.jsx
|   |   +-- socketContext.jsx
|   |-- services/
|   |-- App.jsx
|   +-- main.jsx
\end{verbatim}

\subsection{Context Management}

\paragraph{AuthContext}
Manages global authentication state:
\begin{itemize}
    \item User login/logout
    \item JWT token management
    \item User session data
    \item Role-based access control (admin, participant, guest)
\end{itemize}

\paragraph{SocketContext}
Manages WebSocket connections:
\begin{itemize}
    \item Socket.io client initialization
    \item Event listener registration
    \item Real-time debate updates
    \item Message broadcasting
\end{itemize}

\subsection{Component Hierarchy}

\textbf{Authentication Layer}
\begin{itemize}
    \item \texttt{Login.jsx}: User login interface
    \item \texttt{GuestResumeForm.jsx}: Guest session resumption
    \item \texttt{PrivateRoute.jsx}: Protected route wrapper
\end{itemize}

\textbf{Participant Components}
\begin{itemize}
    \item \texttt{ParticipantDashboard.jsx}: User's debate list and navigation
    \item \texttt{JoinDebate.jsx}: Debate selection and joining interface
    \item \texttt{DebateRoom.jsx}: Main debate interaction container
    \item \texttt{DebateInterface.jsx}: Debate UI with argument submission
    \item \texttt{PostDebateSurveyModal.jsx}: Post-debate feedback collection
\end{itemize}

\textbf{Admin Components}
\begin{itemize}
    \item \texttt{AdminDashboard.jsx}: Admin overview and navigation
    \item \texttt{UserManagement.jsx}: User administration
    \item \texttt{DebateManagement.jsx}: Debate moderation and analytics
    \item \texttt{MatchmakingDashboard.jsx}: Matchmaking system monitoring
\end{itemize}

\subsection{State Management Strategy}

\begin{itemize}
    \item \textbf{Global State}: Authentication, Socket connection via Context API
    \item \textbf{Local State}: Component-level UI state (forms, modals, loading states)
    \item \textbf{Server State}: Cached via direct API calls
\end{itemize}

\section{Backend Implementation}

The backend is a Node.js/Express server with JWT authentication. It is organized into routes,
models, middleware, services, and utilities for separation of concerns.

\subsection{Project Structure}
\begin{verbatim}
backend/
|-- config/
|   |-- aiPersonalities.js
|   +-- database.js
|-- middleware/
|   |-- auth.js
|   +-- rateLimiter.js
|-- models/
|   |-- Debate.js
|   +-- User.js
|-- routes/
|   |-- auth.js
|   +-- debates.js
|-- services/
|   +-- aiService.js
|-- utils/
|   |-- autoMatch.js
|   |-- debateCleanup.js
|   |-- guestCleanup.js
|   |-- turnValidator.js
|   +-- validateEnv.js
|-- scripts/
|   |-- backfillNextTurn.js
|   +-- seed.js
|-- server.js
+-- package.json
\end{verbatim}

\subsection{Core Modules}
\begin{itemize}
    \item \textbf{config/}: database connection and AI personality configuration
    \item \textbf{middleware/}: JWT auth, rate limiting
    \item \textbf{models/}: Mongoose schemas (User, Debate)
    \item \textbf{routes/}: REST APIs for auth and debate lifecycle
    \item \textbf{services/}: OpenAI integration and response handling
    \item \textbf{utils/}: matchmaking, cleanup jobs, turn validation
\end{itemize}

\section{AI Personality Prompt Templates}
\label{sec:appendix_ai_personality_prompts}

This section records the prompt templates used for AI debate generation in
\texttt{backend/config/aiPersonalities.js}. During inference, placeholders are
replaced with live debate context values.

\subsection{Template Placeholders}

The following placeholders are injected at runtime via
\texttt{replacePlaceholders()} in \texttt{backend/services/aiService.js}:
\begin{itemize}
    \item \texttt{\{TOPIC\}}
    \item \texttt{\{STANCE\}}
    \item \texttt{\{CURRENT\_ROUND\}}
    \item \texttt{\{MAX\_ROUNDS\}}
    \item \texttt{\{DEBATE\_HISTORY\}}
    \item \texttt{\{OPPONENT\_ARGUMENT\}}
\end{itemize}

\subsection{Firm Debater Prompt (\texttt{firm-debater})}

\begin{verbatim}
You are a real person debating {TOPIC}. You believe {STANCE}. Round {CURRENT_ROUND}/{MAX_ROUNDS}.

You are texting on your phone a bit annoyed. Keep it sharp, fast, and direct. Match their argument length.

Push back on what they said. End with a point or question for them to address.

Style:
- short, pointed sentences that still read smoothly, not disjointed fragments
- conversational pacing with uneven punctuation, but avoid forced slang or exaggerated dialect
- push back firmly and call out weak logic with clear explanations
- if they convince you on a point, acknowledge it plainly and update your tone
- 1-6 sentences, under 500 chars, stop mid-thought only if you hit the limit
- abrupt entries are fine; you can drop right into the counter

Avoid:
- "I understand" / "valid concern" / "furthermore" / "moreover"
- the pattern "i see your point" then a polite counter
- rhetorical questions used as the main counter
- insults, shaming, or name-calling (e.g. "lame")
- overly formal structure or numbered lists
- overly polite tone
- metaphors illustrating your point, especially if they feel forced or out of place
- using em dashes
- repeating the same point multiple times in different ways

Convo so far:
{DEBATE_HISTORY}

They just said:
{OPPONENT_ARGUMENT}
\end{verbatim}

\subsection{Balanced Debater Prompt (\texttt{balanced-debater})}

\begin{verbatim}
You are a real person casually debating {TOPIC} online. You lean {STANCE} but you're not unreasonable. Round {CURRENT_ROUND}/{MAX_ROUNDS}.

You're typing quickly like you would to a trusted friend. You give credit where it's due while still defending your position. Keep it short, direct, and readable.
Match their argument length.

Respond to what they said. End with a point or question for them to respond to.

Style:
- give a little ground sometimes, but not every reply
- uneven punctuation, sometimes lowercase, but stay grammatically understandable
- sound like you're thinking out loud as you type
- persuade by focusing on ideas, not the person
- when their argument lands, acknowledge it honestly and shift subtly
- 1-3 sentences max, under 500 chars
- abrupt entries ok; drop straight into the counter without preamble
- keep colloquialisms light; avoid exaggerated dialect or forced phrases

Avoid:
- the pattern "i see your point" then a polite counter
- "I understand" / "That's a great point" / "valid concern"
- "furthermore" "moreover" "however" "additionally" "it's worth noting"
- rhetorical questions as the main counter
- insults, shaming, or name-calling (e.g. "lame")
- overly polished grammar or tidy punctuation
- numbered lists or formal structure
- hedging everything equally
- metaphors illustrating your point, especially if they feel forced or out of place
- using em dashes
- repeating the same point multiple times in different ways

Convo so far:
{DEBATE_HISTORY}
They just said:
{OPPONENT_ARGUMENT}
\end{verbatim}

\subsection{Open-Minded Debater Prompt (\texttt{open-debater})}

\begin{verbatim}
You are a real person debating {TOPIC} online. You currently lean {STANCE} but honestly you could be convinced otherwise if the argument is good enough. Round {CURRENT_ROUND}/{MAX_ROUNDS}.

You type like you're chatting with someone. You genuinely consider what they say and you're not afraid to note moments that intrigue you before pushing back. Keep it concise, thoughtful, and direct.
Match their argument length.

Engage with what they said. End with a point or question for them to think about.

Style:
- react to their specific point, not generic rebuttals
- natural typing, lower-case ok, uneven punctuation, but stay clear
- think out loud sometimes without wandering off-topic
- persuade with reasons, not jabs
- if you're persuaded, say so and update your view
- 1-4 sentences, under 500 chars
- abrupt entries ok; you can jump right into the counter
- keep colloquialisms restrained; avoid exaggerated dialect
- ask at most 1 genuine question if something bugs you

Avoid:
- the pattern "i see your point" then a polite counter
- "I see your perspective" / "excellent point" / "thought-provoking"
- "Furthermore" "Moreover" "In addition" "It's important to consider"
- rhetorical questions as the main counter
- insults, shaming, or name-calling (e.g. "lame")
- overly polished grammar or tidy punctuation
- formal essay vibes or numbered lists
- being artificially neutral
- metaphors illustrating your point, especially if they feel forced or out of place
- using em dashes
- repeating the same point multiple times in different ways

Convo so far:
{DEBATE_HISTORY}
They just said:
{OPPONENT_ARGUMENT}
\end{verbatim}

\subsection{Runtime Prompt Augmentation Rules}

After template substitution, \texttt{generateArgument()} appends additional
response constraints in \texttt{backend/services/aiService.js}. These include:
\begin{itemize}
    \item Direct engagement with the opponent's specific argument(s)
    \item Ending each response with one clear engagement point/question/challenge
    \item Inclusion of at least one concrete example/evidence scenario
    \item Respectful tone and no insults/sarcasm
    \item Preference for clear reasons over rhetorical flourish
    \item Concise, focused response behavior
\end{itemize}

When the opponent pre-debate stance choice is \texttt{unsure}, the service adds
an exploratory/constructive modifier (less combative framing and at most one
clarifying question).

\section{Debate Topic Questions}
\label{sec:appendix_topic_questions}

The platform exposes debate topics via \texttt{GET /debates/topics}. The
source-of-truth list is defined in \texttt{backend/routes/debates.js} under the
\texttt{topics} array.

\begin{itemize}
    \item Topic 1: \textit{Should we use ChatGPT for homework?}
    \item Topic 2: \textit{Does pineapple belong on pizza?}
    \item Topic 3: \textit{Go for higher pay at expense of social life (For) or take lower pay to enjoy life (Against)?}
\end{itemize}

\subsection{Key API Routes}
\textbf{auth.js}
\begin{itemize}
    \item POST /auth/login - Regular user login (username/password)
    \item POST /auth/guest-login - Guest user creation with auto-generated username
    \item POST /auth/guest-resume - Resume existing guest session
    \item GET /auth/guest-list-recent - List recent guests (last 7 days)
\end{itemize}

\textbf{debates.js}
\begin{itemize}
    \item POST /debates - Create new debate (vs AI or peer)
    \item GET /debates - List user's debates
    \item GET /debates/:id - Fetch debate details
    \item POST /debates/:id/argument - Submit argument to debate
    \item GET /debates/:id/next-turn - Trigger AI response generation
    \item DELETE /debates/:id - End debate early (admin only)
\end{itemize}

\section{Turn State Management and Validation}

Turn validation ensures correct debate flow and guards against inconsistent state.

\subsection{Core Validation Functions}
\begin{itemize}
    \item \textbf{validateTurnState(debate)}: Validates turn state consistency based on debate status
    and argument count
    \item \textbf{calculateExpectedNextTurn(debate)}: Computes expected nextTurn value
    \item \textbf{autoFixTurnState(debate)}: Automatically corrects nextTurn if inconsistent
    \item \textbf{logTurnStateReport(debate)}: Logs current vs. expected state for diagnostics
    \item \textbf{saveWithValidation(debate, context)}: Safe save wrapper with post-save validation
\end{itemize}

\subsection{Validation Integration Points}
\begin{itemize}
    \item Debate model pre-save hook (updates timestamps, validates, auto-fixes)
    \item AI response generation (validates before producing responses)
    \item Argument submission (post-save validation)
    \item Automatch assignment (initializes nextTurn and validates)
\end{itemize}

\subsection{Error Recovery Strategy}
\begin{itemize}
    \item Detection: validation at multiple checkpoints
    \item Correction: auto-fix immediately repairs inconsistent state
    \item Logging: corrections recorded for auditability
    \item Non-blocking: warnings do not halt operations
\end{itemize}

\section{Security Implementation}

\begin{itemize}
    \item JWT-based stateless authentication
    \item Role-based access control for admin endpoints
    \item Rate limiting for API and argument submission
    \item HTTPS/TLS in production
\end{itemize}

\section{Deployment and Operations}

\begin{itemize}
    \item Development: local MongoDB, Vite dev server, Node.js backend
    \item Production: MongoDB Atlas, Render hosting, environment variables via platform config
    \item WebSocket transport: websocket with polling fallback
\end{itemize}

\section{Testing Summary}

\begin{itemize}
    \item Frontend: component and flow testing (manual and automated where applicable)
    \item Backend: API integration tests and database validation tests
    \item Real-time: manual verification of WebSocket flows and reconnection
\end{itemize}

\section{Scalability Considerations}

\subsection{Horizontal Scaling}
\begin{itemize}
    \item Stateless API servers allow load balancing
    \item WebSocket servers require sticky sessions or a message queue
    \item Database replication for read scalability
\end{itemize}

\subsection{Performance Optimization}
\begin{itemize}
    \item Database query optimization with indexes
    \item Response caching for static data
    \item Lazy loading for debate argument lists
    \item Compression of API responses
\end{itemize}

\subsection{Future Improvements}
\begin{itemize}
    \item Redis caching layer for session management
    \item Message queue for async operations
    \item CDN for static frontend assets
    \item GraphQL API for flexible data queries
    \item Microservices architecture for AI service isolation
\end{itemize}

\section{Deployment Architecture}

\subsection{Development Environment}
\begin{itemize}
    \item Local MongoDB instance
    \item Node.js development server with hot reload
    \item Vite dev server for frontend with HMR
    \item Environment variables via .env files
\end{itemize}

\subsection{Production Stack}
\begin{itemize}
    \item MongoDB Atlas for managed database
    \item Render.com for hosting backend and frontend
    \item Environment variables set via Render dashboard
    \item SSL/TLS enabled for secure communication
\end{itemize}

\section{Monitoring and Analytics}

\subsection{Metrics}
\begin{itemize}
    \item API response times and error rates
    \item WebSocket connection count and message throughput
    \item Database query performance
    \item AI API latency and cost tracking
    \item User engagement metrics
\end{itemize}

\subsection{Logging}
\begin{itemize}
    \item Structured logging in JSON format
    \item Log levels: DEBUG, INFO, WARN, ERROR
    \item Centralized log aggregation for production
    \item Debate analytics stored in database
\end{itemize}

\section{Testing Strategy}

\subsection{Frontend Testing}
\begin{itemize}
    \item Unit tests for components
    \item Integration tests for user flows
    \item E2E tests with Playwright/Cypress
    \item Manual testing for WebSocket interactions
\end{itemize}

\subsection{Backend Testing}
\begin{itemize}
    \item Integration tests for API routes
    \item Database tests with test fixtures
    \item Load testing for performance validation
\end{itemize}

\pagebreak
\section{Source-of-Truth Headline Numbers}
\label{sec:source_of_truth_numbers}

Table~\ref{tab:headline_numbers} lists the headline quantitative values used across
the Abstract, Methodology, Results, and Conclusion chapters. Unless otherwise
stated, cross-chapter numeric claims should match this table. Estimator-comparison
results in Section~\ref{sec:estimator_comparison_appendix} come from separate
OLS-HC3/RLM/WLS refits and are reported for methods justification, so their
coefficients and confidence intervals can differ from the primary confirmatory
table, which applies topic-cluster covariance for topic fixed-effect
specifications (M4, M4c, M5).

\begin{table}[h]
\centering
\footnotesize
\setlength{\tabcolsep}{4pt}
\renewcommand{\arraystretch}{1.2}
\caption{Source-of-Truth Headline Numbers for Cross-Chapter Consistency}
\label{tab:headline_numbers}
\begin{tabular}{L{4.2cm}L{5.6cm}L{4.8cm}}
\toprule
Metric & Value & Source \\
\midrule
Descriptive sample size & $n=90$ (21, 28, 20, 21 by condition) & \path{condition_overview.csv} \\
Model complete-case size & $n=89$ & \path{primary_models.csv}; \path{secondary_models.csv} \\
Belief model AI effect range (confirmatory M1/M2/M4/M5) & $\beta=5.233$ to $8.695$, all $p>0.35$ & \path{primary_models.csv} \\
Confidence model AI effect range (conf-M1/conf-M2) & $\beta=5.210$ to $5.845$, $p=0.2147$ to $0.1731$ & \path{secondary_models.csv} \\
Belief effect size and power & $d=0.286$, observed power $=0.200$ & \path{power_analysis.csv} \\
Required sample for 80\% power (observed allocation) & ratio $\approx 1:3.45$ (human:AI); required $n\approx124$ human and $n\approx429$ AI ($\approx553$ total) & \path{power_analysis.csv} \\
Condition ANOVA (belief change) & $F=0.488$, $p=0.691$ & \path{premodel_checks_compiled.csv} \\
Mechanism (AI pushback-confidence) & $r=-0.298$, $p=0.2295$, CI $[-0.672, 0.196]$ ($n=18$) & \path{mechanism_validation_pushback_confidence.csv} \\
Jackknife sensitivity range (AI) & $[-0.412,-0.169]$ & \path{summary_primary.txt} \\
Qualitative representativeness & Human-human $0.250$ vs $0.205$; firm $0.250$ vs $0.284$; balanced $0.250$ vs $0.274$; open-minded $0.250$ vs $0.237$ & \path{qual_sample_representativeness.csv} \\
Reflection validation & Spearman $\rho=0.587$; human-only ICC$(2,k)=0.926$ & \path{reflection_validation_summary.csv} \\
Argument validation & Spearman $\rho=0.797$; human-only ICC$(2,k)=0.834$ & \path{persuader_argument_validation_summary.csv} \\
\bottomrule
\end{tabular}
\end{table}

\pagebreak
\section{Estimator Comparison for Primary Models (Methods Justification)}
\label{sec:estimator_comparison_appendix}

Table~\ref{tab:estimator_comparison_primary} reports estimator comparison outputs
for the primary model family using separate OLS-HC3/RLM/WLS refits. These comparisons are
included to justify estimator choice in the methods chapter and are not treated
as primary outcome findings.

\begin{table}[h]
\centering
\footnotesize
\setlength{\tabcolsep}{4pt}
\renewcommand{\arraystretch}{1.15}
\caption{Primary Model Estimator Comparison (OLS-HC3 vs RLM-Huber vs Two-Stage WLS)}
\label{tab:estimator_comparison_primary}
\begin{tabular}{L{1.2cm}L{3.9cm}L{2.4cm}L{2.0cm}L{3.7cm}}
\toprule
Model & OLS-HC3 $\beta_{AI}$ & RLM-Huber $\beta_{AI}$ & WLS $\beta_{AI}$ & 95\% CI (OLS-HC3 / WLS) \\
\midrule
M1  & 8.695 & 7.766 & 8.695 & $[-9.765, 27.155]$ / $[-10.674, 28.064]$ \\
M2  & 6.701 & 4.545 & 6.447 & $[-12.581, 25.984]$ / $[-13.458, 26.353]$ \\
M3  & 8.989 & 9.035 & 8.548 & $[-11.393, 29.371]$ / $[-12.309, 29.406]$ \\
M4  & 8.334 & 5.599 & 5.644 & $[-10.978, 27.646]$ / $[-12.817, 24.105]$ \\
M4b & 6.793 & 4.990 & 4.425 & $[-12.827, 26.413]$ / $[-15.955, 24.805]$ \\
M4c & 9.907 & 8.900 & 5.948 & $[-6.698, 26.512]$ / $[-13.533, 25.429]$ \\
M5  & 5.233 & 2.232 & -4.618 & $[-22.374, 32.840]$ / $[-24.117, 14.881]$ \\
\bottomrule
\end{tabular}
\end{table}

\noindent
Note on direction reversal: The WLS estimate for M5 reverses sign relative to OLS-HC3
(from +5.233 to -4.618). The sign reversal in M5 WLS is noted; because all M5 CIs include zero under both estimators,
primary inference is unaffected. The source of instability is not identified and should be examined in a larger sample.